\par\medskip
\noindent
14. \underline{Relaciones con la teoría trascendente: el caso
analítico rígido}. \underline{El homomorfismo canónico}
$\underline M_K\longrightarrow A_K^{\natural\mathrm o}$.

\medskip
\noindent
14.1. Suponemos aquí conocidos los fundamentos de la teoría de los
espacios analíticos rígidos, para los cuales pueden consultarse [12],
[13] y [34].

Sea $S$ un rasgo completo y sea $A_K$ un esquema abeliano sobre su
cuerpo de fracciones. En 8.1, mediante descenso a partir de la extensión
de Raynaud en el caso semiestable, construimos una extensión
$G_K=A_K^{\natural\mathrm o}$ (8.1.1)

\begin{equation}\tag{14.1.1}
 0\longrightarrow T_K\longrightarrow G_K\longrightarrow B_K
 \longrightarrow0
\end{equation}

\noindent
de un esquema abeliano $B_K$ por un toro $T_K$. Considerando el esquema
abeliano dual $A'_K$, se tiene asimismo una extensión
$G'_K=A'{}^{\natural\mathrm o}_K$

\begin{equation}\tag{14.1.2}
 0\longrightarrow T'_K\longrightarrow G'_K\longrightarrow B'_K
 \longrightarrow0\qquad ,
\end{equation}

\noindent
donde, además, $B'_K$ se identifica con el esquema abeliano dual de
$B_K$ (8.2). Como en (8.1.8), sean $\underline M_K$ y
$\underline M'_K$ las redes sobre $K$ duales de $T'_K$ y $T_K$,
respectivamente:

\begin{equation}\tag{14.1.3}
 \underline M_K=D(T'_K),\qquad \underline M'_K=D(T_K)\qquad .
\end{equation}

\noindent
Así, las estructuras de extensión (14.1.2) y (14.1.1) pueden
interpretarse respectivamente (VIII 3.7) como equivalentes a los datos
de homomorfismos

\begin{equation}\tag{14.1.4}
 \underline M_K\xrightarrow{\ p_K\ }B_K,\qquad
 \underline M'_K\xrightarrow{\ p'_K\ }B'_K\qquad .
\end{equation}

Sea $W_K^{\natural}$ la biextensión de $(B_K,B'_K)$ por
$\mathbb G_{mK}$ que define la dualidad canónica entre estos esquemas
abelianos (8.2), y sea $E_K$ la biextensión de
$(\underline M_K,\underline M'_K)$ por $\mathbb G_{mK}$ obtenida al
traer $W_K^{\natural}$ por los homomorfismos (14.1.4). Se comprueba
inmediatamente que las trivializaciones de $E_K$ corresponden
canónica y biyectivamente a homomorfismos

\begin{equation}\tag{14.1.5}
 i:\underline M_K\longrightarrow G_K
\end{equation}

\noindent
que levantan $p_K:\underline M_K\longrightarrow B_K$, o,
simétricamente, a homomorfismos

\begin{equation}\tag{14.1.6}
 i':\underline M'_K\longrightarrow G'_K
\end{equation}

\noindent
que levantan $p'_K:\underline M'_K\longrightarrow B'_K$.

Sentado esto, escribamos un superíndice ${}^{\mathrm{an}}$ para los
grupos analíticos rígidos sobre $K$ asociados a los grupos algebraicos
considerados [13]. En [24] se define una sucesión exacta canónica

\begin{equation}\tag{14.1.7}
 0\longrightarrow\underline M_K^{\mathrm{an}}
 \xrightarrow{\ i^{\mathrm{an}}\ }G_K^{\mathrm{an}}
 \xrightarrow{\ e\ }A_K^{\mathrm{an}}\longrightarrow0\qquad .
\end{equation}

\noindent
Aquí $i^{\mathrm{an}}$, que también puede interpretarse como un
homomorfismo $i:\underline M_K\longrightarrow G_K$ de esquemas en
grupos ordinarios, \underline{levanta el homomorfismo} $p$ de (14.1.4),
es decir, es de la forma considerada en (14.1.5). Esta sucesión exacta
es el análogo analítico rígido de la sucesión exacta (13.3.4). Depende
funtorialmente de $A_K$, y su formación es compatible con todo cambio
de rasgos $S'\longrightarrow S$.

El último hecho muestra también que, para definir y estudiar (14.1.7),
el descenso finito nos reduce al caso en que $A_K$ posee reducción
semiestable. En este caso, los esquemas en grupos
$T_K,G_K,B_K,\underline M_K,\ldots$ son las fibras genéricas de
esquemas en grupos $T,G,B,M,\ldots$ sobre $S$, siendo $G$ una extensión
del esquema abeliano $B$ por el toro $T$ (párrafo 7). Podemos considerar
entonces el completado formal $\hat G$ de $G$ a lo largo de su fibra
especial, y por construcción existe un isomorfismo canónico (7.2.6)

\begin{equation}\tag{14.1.8}
 \hat G\xrightarrow{\ \sim\ }\widehat{A^{\mathrm o}}
 \qquad .
\end{equation}

\noindent
Al pasar a las «fibras genéricas», que son grupos analíticos rígidos
sobre $K$, se obtiene un isomorfismo canónico

\begin{equation}\tag{14.1.9}
 \hat G^{\mathrm{an}}\xrightarrow{\ \sim\ }
 \widehat{A^{\mathrm o}}\,{}^{\mathrm{an}}
 \qquad .
\end{equation}

\noindent
Por otra parte, existen homomorfismos canónicos (inmersiones abiertas
analíticas rígidas, pero no isomorfismos en general)

\begin{equation}\tag{14.1.10}
 \hat G^{\mathrm{an}}\lhook\joinrel\longrightarrow G_K^{\mathrm{an}},
 \qquad
 \widehat{A^{\mathrm o}}\,{}^{\mathrm{an}}
 \lhook\joinrel\longrightarrow A_K^{\mathrm{an}}
 \qquad .
\end{equation}

\noindent
El homomorfismo $e$ de (14.1.7) se caracteriza de manera única por la
conmutatividad del diagrama

\begin{equation}\tag{14.1.11}
\begin{tikzcd}
 \hat G^{\mathrm{an}} \arrow[r] \arrow[d]
   & \widehat{A^{\mathrm o}}\,{}^{\mathrm{an}} \arrow[d] \\
 G_K^{\mathrm{an}} \arrow[r] & A_K^{\mathrm{an}}
\end{tikzcd}
\end{equation}

\noindent
donde las flechas sin nombre son las de (14.1.9) y (14.1.10).

\medskip
\noindent
\underline{Observaciones 14.1.12.}

\smallskip
\noindent
a) Obsérvese que la unicidad de $e$ que hace conmutativo el diagrama
precedente es inmediata, pues el espacio analítico rígido
$G_K^{\mathrm{an}}$ es conexo. En cambio, la existencia de $e$ no es en
absoluto trivial, y sobre este punto nos limitamos a remitir a la nota
citada de M.~Raynaud. Además, la definición de $i^{\mathrm{an}}$ dada
en loc. cit. es asimismo de naturaleza trascendente, aunque puede
interpretarse como un objeto de naturaleza esencialmente algebraica
(14.1.5), una vez obtenida la construcción del propio $G$, y por tanto
de $G_K$, mediante las construcciones de la \underline{geometría formal}
del párrafo 7. No conozco ninguna descripción puramente
algebrogeométrica de este homomorfismo, pero más adelante (14.4)
propondré una caracterización «aritmética», obtenida por reducción a una
situación de tipo finito sobre el anillo $\mathbb Z$. Así pues, desde el
punto de vista de la geometría algebraica, el homomorfismo (14.1.5)
parece ser una estructura más oculta que las encontradas en los
párrafos precedentes.

\smallskip
\noindent
b) Como se explicó arriba, el homomorfismo $i$ está asociado
canónicamente a una trivialización bien determinada de la biextensión
$E_K$ de $(\underline M_K,\underline M'_K)$, o, equivalentemente, a un
homomorfismo bien determinado (14.1.6) que levanta $p'_K$. Se demuestra
entonces que este último es precisamente el homomorfismo que aparece en
la sucesión exacta

\begin{equation}\tag{14.1.13}
 0\longrightarrow\underline M'{}^{\mathrm{an}}_K
 \xrightarrow{\ i'{}^{\mathrm{an}}\ }G'{}^{\mathrm{an}}_K
 \longrightarrow A'{}^{\mathrm{an}}_K\longrightarrow0
\end{equation}

\noindent
análoga a (14.1.7), asociada al esquema abeliano $A'_K$ dual de $A_K$
(¿a menos que sea el opuesto???).
